%==================================================================

\begin{thm}[главное утверждение лекции]\label{thm:hom}
Пусть $\Shv{F}$ --- предпучок на $X$, а $\Shv{G}$ --- \emph{пучок} на
$X$. Тогда отображение
\begin{equation}\label{eq:Phi}
 \Phi\colon\Hom(\Shv{F}^{\#},\Shv{G})\longrightarrow
 \Hom(\Shv{F},\Shv{G}),\qquad \varphi\longmapsto\varphi\circ i,
\end{equation}
где $i\colon\Shv{F}\to\Shv{F}^{\#}$ --- каноническое, является
изоморфизмом. Более того, изоморфизм естествен: коммутативны
квадраты диаграмм~\ref{dgm:natG} и~\ref{dgm:natF}.
\end{thm}

Заметим направление: изоморфизм идёт \emph{из} морфизмов в пучок
$\Shv{F}^{\#}$ \emph{в} морфизмы в произвольный предпучок. Гомоморфизм
пучка определяется своими значениями на сечениях, а сечений
$\Shv{F}^{\#}(U)$ <<достаточно много>>: именно это и доказывается.

\subsection{Шаг 1: инъективность $\Phi$}
Для $\varphi\colon\Shv{F}^{\#}\to\Shv{G}$ положим
$\varphi':=\varphi\circ i\colon\Shv{F}\to\Shv{G}$ --- диаграмма
\ref{dgm:univ}.

\begin{diagram}{Диаграмма универсальности: $\varphi'=\varphi\circ i$
 --- сжатие $\varphi$ по каноническому $i$.\label{dgm:univ}}
\begin{tikzcd}[cd style,column sep=huge,row sep=huge]
 \Shv{F} \ar[r,"i"] \ar[dr,"\varphi'"'] & \Shv{F}^{\#} \ar[d,"\varphi"'] \\
 & \Shv{G}
\end{tikzcd}
\end{diagram}

Пояснение к диаграмме~\ref{dgm:univ}. Треугольник читается как
<<сначала сжать по $i$, затем применить $\varphi$>>: диагональ
$\varphi'=\varphi\circ i$ --- это и есть $\Phi(\varphi)$, переход
морфизма пучка в морфизм предпучков. Коммутативность треугольника
--- просто определение $\varphi'$; всё содержание утверждения
в том, что $\varphi'$ корректно определён и что из
$\varphi\circ i=\psi\circ i$ следует $\varphi=\psi$ ---
доказывается сразу ниже.

Если $\varphi\circ i=\psi\circ i$, то $\varphi=\psi$: действительное
значение $\varphi$ на произвольном сечении $\tilde s\in\Shv{F}^{\#}(U)$
уже определено, потому что $\tilde s$ локально является образом
некоторого $s\in\Shv{F}(V)$, а $\varphi$ --- гомоморфизм \emph{пучка}
$\Shv{G}$, поэтому его поведение локально, а значит и везде,
определяется значениями $i(s)$. Разделение точек $\Phi(\varphi)$ ---
это инъективность.

\subsection{Шаг 2: сюръективность $\Phi$}
\subsubsection*{Построение $\beta$}

Пусть дан $\alpha\colon\Shv{F}\to\Shv{G}$ --- гомоморфизм предпучков.
Надо построить $\beta\colon\Shv{F}^{\#}\to\Shv{G}$ с
$\beta\circ i=\alpha$.

Начнём с <<наивной>> формулы. Для $U\subset X$ и
$s\in\Shv{F}(U)$ положим
\begin{equation}\label{eq:beta-naive}
 \beta\bigl(i(s)\bigr):=\alpha_U(s)\in\Shv{G}(U).
\end{equation}
Это хорошо согласуется с ограничениями, потому что $\alpha$ ---
гомоморфизм предпучков: $\beta(i(s))|_{V}=\alpha_V(s|_V)$.

\subsubsection*{Проблема единственности}

Возникает вопрос: а что, если $s_1\neq s_2$, но $i(s_1)=i(s_2)$? Тогда
по~\eqref{eq:beta-naive} у $\beta$ получится два разных значения
$g_1=\alpha(s_1)$ и $g_2=\alpha(s_2)$, и отображение не будет
определено. Схема на листе~11: $s_1\to\tilde s\leftarrow s_2$.
Покажем, что в таком случае $g_1=g_2$.

\begin{lem}[локальное равенство]\label{lem:local}
Пусть $s,t\in\Shv{F}(U)$ и $i(s)=i(t)$ в $\Shv{F}^{\#}(U)$. Тогда
$\alpha_U(s)=\alpha_U(t)$ в $\Shv{G}(U)$.
\end{lem}

\begin{proof}
По условию ростки совпадают: $s_x=t_x$ для всех $x\in U$. По
определению прямого предела~\eqref{eq:stalk} это означает: для каждого
$x\in U$ существует открытая окрестность $V_x\subset U$ такая, что
\begin{equation}\label{eq:same}
 \res^{V_x}_{U}(s)=\res^{V_x}_{U}(t),
 \qquad\text{т.\,е.}\qquad s|_{V_x}=t|_{V_x}.
\end{equation}
Семейство $\{V_x\}_{x\in U}$ покрывает $U$. Применив к~\eqref{eq:same}
гомоморфизм предпучка $\alpha$ (точнее, его компоненту на $V_x$), получим
\[
 \alpha_{V_x}\bigl(s|_{V_x}\bigr)=\alpha_{V_x}\bigl(t|_{V_x}\bigr),
\]
потому что $\alpha$ переводит равные элементы в равные. С другой
стороны, коммутативность диаграммы ограничений для $\alpha$
(диаграмма~\ref{dgm:natur}) даёт
\[
 \alpha_{V_x}\bigl(s|_{V_x}\bigr)=\bigl(\alpha_U(s)\bigr)|_{V_x},
 \qquad
 \alpha_{V_x}\bigl(t|_{V_x}\bigr)=\bigl(\alpha_U(t)\bigr)|_{V_x}.
\]
Значит $\alpha_U(s)$ и $\alpha_U(t)$ совпадают на каждом $V_x$:
$\bigl(\alpha_U(s)\bigr)|_{V_x}=\bigl(\alpha_U(t)\bigr)|_{V_x}$.
Осталось применить аксиому \emph{(F1) пучка $\Shv{G}$} к покрытию
$U=\bigcup_xV_x$: два глобальных сечения, совпадающие на каждом куске
покрытия, равны. Итак, $\alpha_U(s)=\alpha_U(t)$.
\end{proof}

\begin{diagram}{Натурность гомоморфизма $\alpha\colon\Shv{F}\to\Shv{G}$:
 квадраты коммутативны для любой пары $W\subset W'$, в частности для
 $V_x\subset U$. Слева направо --- ограничения, сверху вниз ---
 компоненты $\alpha$.\label{dgm:natur}}
\begin{tikzcd}[cd style]
 \Shv{F}(U) \ar[r,"\res^{V_x}_{U}"] \ar[d,"\alpha_U"'] &
 \Shv{F}(V_x) \ar[d,"\alpha_{V_x}"] \\
 \Shv{G}(U) \ar[r,"\res^{V_x}_{U}"'] & \Shv{G}(V_x)
\end{tikzcd}
\end{diagram}

Пояснение к диаграмме~\ref{dgm:natur}. Это квадрат натуральности
гомоморфизма предпучков $\alpha$: два маршрута из $\Shv{F}(U)$
в $\Shv{G}(V_x)$ --- «сначала ограничить, потом применить
$\alpha$» (вверх-вправо-вниз) и «сначала применить $\alpha$ на
$U$, потом ограничить» (влево-вниз-направо) --- дают одно и то
же. Это ровно условие $\alpha_U(s)|_{V_x}=\alpha_{V_x}(s|_{V_x})$;
в доказательстве квадрат применяется к паре $V_x\subset U$.

Лемма~\ref{lem:local} означает: формула~\eqref{eq:beta-naive} не зависит
от того, каким из $s$ с $\tilde s=i(s)$ мы воспользуемся.

\subsubsection*{Определение $\beta$ на произвольном сечении}

Теперь возьмём произвольное $\tilde s\in\Shv{F}^{\#}(U)$, а не только
вида $i(s)$. Почему оно локально такого вида? Потому что непрерывность:
для точки $x\in U$ росток $\tilde s(x)\in\Shv{F}_x$ --- это класс
некоторой пары $(V,s)$, $s\in\Shv{F}(V)$; множество точек, в которых
$\tilde s$ совпадает с $i(s)$, открыто (см.~шаг~1 доказательства
(F1) в \S\,\ref{sec:axioms}---там показано, что множество точек
совпадения двух непрерывных сечений открыто). Эти открытые множества
покрывают $U$, значит $U=\bigcup_\alpha W_\alpha$ и
$\tilde s|_{W_\alpha}=i(s_\alpha)$ для некоторых
$s_\alpha\in\Shv{F}(W_\alpha)$.

Положим
\[
 t_\alpha:=\alpha_{W_\alpha}(s_\alpha)\in\Shv{G}(W_\alpha).
\]
На пересечении $W_\alpha\cap W_\beta$ выполняется
$i(s_\alpha)|=i(s_\beta)|$, т.\,е. локальное равенство сечений
$\Shv{F}$; по лемме~\ref{lem:local} (вернее, по её локальной версии)
$\alpha(s_\alpha)$ и $\alpha(s_\beta)$ совпадают на
$W_\alpha\cap W_\beta$. Итак, семейство $(t_\alpha)$ --- согласованное,
и аксиома \emph{(F2) пучка $\Shv{G}$} даёт единственное
$t\in\Shv{G}(U)$ с $t|_{W_\alpha}=t_\alpha$. Полагаем
\begin{equation}\label{eq:beta}
 \beta(\tilde s):=t.
\end{equation}

Независимость от выбора покрытия: если есть другое покрытие
$U=\bigcup_\beta W'_\beta$, то на пересечениях $W_\alpha\cap W'_\beta$
оба определения совпадают (там $\tilde s$ --- образ одного и того же
сечения), и снова применяется (F2), либо (F1) --- глобальное сечение
$\Shv{G}(U)$ определяется по значениям на покрытии однозначно.

Итого: $\beta\colon\Shv{F}^{\#}\to\Shv{G}$ --- гомоморфизм пучков,
и $\beta\circ i=\alpha$ по построению~\eqref{eq:beta}. Значит
$\Phi(\beta)=\alpha$, т.\,е. $\Phi$ --- сюръективно.

\subsubsection*{Схема склейки}

Глобальное сечение $t$ из~\eqref{eq:beta} собирается из локальных
$t_\alpha$ ровно так, как показано на диаграмме~\ref{dgm:glue}: значения
на $W_\alpha$ и $W_\beta$ дают одно и то же на пересечении, и (F2)
<<склеивает>> их в $t\in\Shv{G}(U)$.

\begin{diagram}{Склейка: из локальных сечений $t_x\in\Shv{G}(V_x)$,
 согласованных на $V_x\cap V_y$, аксиома (F2) пучка $\Shv{G}$ даёт
 единственное $t\in\Shv{G}(U)$.\label{dgm:glue}}
\begin{tikzcd}[cd style]
 & \Shv{G}(U) \ar[dl,"\res^{V_x}_{U}"'] \ar[dr,"\res^{V_y}_{U}"] & \\
 \Shv{G}(V_x) \ar[dr,"\res"'] & & \Shv{G}(V_y) \ar[dl,"\res"] \\
  & \Shv{G}(V_x\cap V_y) &
\end{tikzcd}
\end{diagram}

Пояснение к диаграмме~\ref{dgm:glue}. Ромб показывает действие
аксиомы (F2) в обе стороны. Вниз: глобальное сечение
$t\in\Shv{G}(U)$ ограничивается до $V_x$ и $V_y$, и оба
полученных локальных значения совпадают на пересечении
$V_x\cap V_y$ (две нижние стрелки сходятся к одному элементу).
Вверх --- сама склейка: если заданы согласованные локальные
значения $t_x$ и $t_y$, то (F2) даёт единственное
$t\in\Shv{G}(U)$, из которого они получаются ограничением.

\sectionsrc{photo\_3---photo\_4 (стр.~11---12):
 проблема $s_1\neq s_2$, $i(s_1)=i(s_2)$; выбор $V_x$ <<по определению
 прямого предела>>; натурная диаграмма; вывод
 <<$\forall x,y\in U\ \exists V_x,V_y$ и $g_x\in\Shv{G}(V_x)$,
 $g_y\in\Shv{G}(V_y)$ с $g_x|_{V_x\cap V_y}=g_y|_{V_x\cap V_y}$;
 применяем (F2) к пучку $\Shv{G}$ и получаем
 $g\in\Shv{G}(U)$, определённый однозначно $\Rightarrow$ отображение
 $\beta$ задано корректно».}

\subsection{Шаг 3: естественность}
Для гомоморфизма пучков $f\colon\Shv{G}_1\to\Shv{G}_2$ и для
гомоморфизма предпучков $g\colon\Shv{F}_1\to\Shv{F}_2$ коммутативны
квадраты диаграмм~\ref{dgm:natG} и~\ref{dgm:natF}.

\begin{diagram}{Естественность $\Phi$ по $\Shv{G}$: сжатие по $i$
 совместно с композицией с $f$.\label{dgm:natG}}
\begin{tikzcd}[cd style]
 \Hom(\Shv{F},\Shv{G}_1) \ar[r,"f\circ-"] \ar[d,"\Phi"'] &
 \Hom(\Shv{F},\Shv{G}_2) \ar[d,"\Phi"] \\
 \Hom(\Shv{F}^{\#},\Shv{G}_1) \ar[r,"f\circ-"'] &
 \Hom(\Shv{F}^{\#},\Shv{G}_2)
\end{tikzcd}
\end{diagram}

Пояснение к диаграмме~\ref{dgm:natG}. Квадрат естественности
$\Phi$ по пучку: верхняя стрелка --- домножение морфизма
предпучков на $f\colon\Shv{G}_1\to\Shv{G}_2$, нижняя --- то же
для морфизмов в $\Shv{G}_2$, вертикали --- <<сжатие>> $\Phi$
по каноническому $i$. Коммутативность означает:
$f\circ(\varphi\circ i)=(f\circ\varphi)\circ i$, т.\,е.\ неважно,
сжимать до или после домножения на $f$.

\begin{diagram}{Естественность $\Phi$ по $\Shv{F}$: сжатие по $i_2$
 совместно с композиции с $g$.\label{dgm:natF}}
\begin{tikzcd}[cd style]
 \Hom(\Shv{F}_1,\Shv{G}) \ar[r,"-\circ i_2"] \ar[d,"\Phi"'] &
 \Hom(\Shv{F}_2,\Shv{G}) \ar[d,"\Phi"] \\
 \Hom(\Shv{F}_1^{\#},\Shv{G}) \ar[r,"-\circ i_2"'] &
 \Hom(\Shv{F}_2^{\#},\Shv{G})
\end{tikzcd}
\end{diagram}

Пояснение к диаграмме~\ref{dgm:natF}. Зеркальная ситуация:
стрелки $-\circ i_2$ домножают морфизм предпучков
$g\colon\Shv{F}_1\to\Shv{F}_2$, вертикали --- снова $\Phi$.
Коммутативность --- это $(\varphi\circ g)\circ i_1
=(\varphi\circ i_2)\circ g$: сжимать сначала $g$, потом $\varphi$,
или наоборот, --- одинаково. Вместе две диаграммы
(\ref{dgm:natG} и \ref{dgm:natF}) и означают, что изоморфизм
$\Phi$ естественен по обоим аргументам.

Здесь $i_2\colon\Shv{F}_2\to\Shv{F}_2^{\#}$ --- каноническое
отображение, а стрелки $-\circ i_2$ индуцированы гомоморфизмом
предпучков $g$. Проверка прямая: $\Phi(f\circ\varphi)=
f\circ\varphi\circ i=f\circ\Phi(\varphi)$ и
$\Phi(\varphi\circ g)=(\varphi\circ g)\circ i_1=(\varphi\circ i_2)\circ
g=\Phi(\varphi)\circ g$.

\begin{rem}[что это значит категорно]
Если задан функтор $\PSh(X)\to\Sh(X)$, отсылающий предпучок к пучку,
порождённому им, то изоморфизм теоремы~\ref{thm:hom} ---
это \emph{свойство универсальности} этого функтора: включение
$\Sh(X)\hookrightarrow\PSh(X)$ имеет левый сопряжённый
$\Shv{F}\mapsto\Shv{F}^{\#}$. Коммутативность диаграмм
\ref{dgm:natG}, \ref{dgm:natF} --- naturality этого построения;
именно она позволяет говорить, что <<прокачивание>> задано
\emph{функторно}, а не произвольным выбором.
\end{rem}

\begin{bookbox}
\booksrc{Годеман, гл.~I, п.~1.7, с.~24---25: категории, ковариантные и
 контравариантные функторы, естественное преобразование (система
 $T(X)$, которая делает диаграмму с $F(f),G(f)$ коммутативной); пример:
 морфизм $t\colon A\to B$ индуцирует морфизм функторов
 $\Hom(B,X)\to\Hom(A,X)$ передачей $u\mapsto u\circ t$.
 Годеман, гл.~I, п.~1.9, с.~30---31: предпучок --- контравариантный
 функтор категории открытых множеств; категория предпучков абелевых
 групп --- абелева категория, функтор $\Shv{F}\mapsto\Shv{F}(U)$ точен,
 поэтому точность последовательности предпучков эквивалентна
 поэлементной точности.
\par\smallskip
 Манин, \S\,1.16.5---1.16.6, с.~103: предпучок как контравариантный
 функтор $\mathrm{Top}_E\to\mathrm{Sets}$; определение функторного
 морфизма с коммутативной диаграммой; категория
 $\mathrm{Funct}(C,D)$.
}
\end{bookbox}

%==================================================================
